\documentclass[14pt]{article}
\usepackage[utf8]{inputenc}	% Para caracteres en español
\usepackage{amsmath,amsthm,amsfonts,amssymb,amscd}
\usepackage{multirow,booktabs}
\usepackage[table]{xcolor}
\usepackage{fullpage}
\usepackage{lastpage}
\usepackage{enumitem}
\usepackage{fancyhdr}
\usepackage{mathrsfs}
\usepackage{wrapfig}
\usepackage{setspace}
\usepackage{calc}
\usepackage{multicol}
\usepackage{cancel}
\usepackage[title]{appendix}
\usepackage[retainorgcmds]{IEEEtrantools}
\usepackage[margin=3cm]{geometry}
\usepackage{amsmath, amsthm}
\newlength{\tabcont}
\setlength{\parindent}{0.0in}
\setlength{\parskip}{0.05in}
\usepackage{empheq}
\usepackage{framed}
\usepackage[most]{tcolorbox}
\usepackage{xcolor}
\colorlet{shadecolor}{orange!15}
\parindent 0in
\parskip 12pt
\geometry{margin=1in, headsep=0.25in}
\newtheorem{theorem}{Theorem}
\newtheorem{lemma}[theorem]{Lemma}
\newtheorem*{remark}{Remark}
\newtheorem{corollary}[theorem]{Corollary}
\newtheorem{proposition}[theorem]{Proposition}
\theoremstyle{definition}
\newtheorem*{definition}{Definition}
\newtheorem{example}[theorem]{Example}
\theoremstyle{remark}
\newtheorem{notation}{Notation}
%\AtAppendix{\counterwithin{theorem}{section}}
\setlength{\parskip}{1em}
\DeclareMathOperator{\sign}{sign}
\renewcommand{\thetheorem}{\arabic{section}.\arabic{theorem}}
\renewcommand{\theequation}{\arabic{section}.\arabic{equation}}
\usepackage[numbib,notlof,notlot,nottoc]{tocbibind}
\numberwithin{equation}{section}
\numberwithin{theorem}{section}
\usepackage{xeCJK}
    \setCJKmainfont[AutoFakeBold=3]{AR PL UKai TW}
    \XeTeXlinebreaklocale "zh"             
    \XeTeXlinebreakskip = 0pt plus 1pt


\begin{document}
\setcounter{section}{0}
\title{Sobolev Embedding Theorem}

\thispagestyle{empty}
\begin{center}
{\LARGE \bf Sobolev Embedding Theorem}
\end{center}
\section{Sobolev space $W^{1,p}(I)$}
\begin{definition} The \emph{Sobolev space} $W^{1,p}(I)$ is defined by
\begin{equation}\label{Sineq}
W^{1,p}(I)=\{u\in L^p(I): there\ exists\ g\in L^p(I)\ such\ that \int_I\,u\,\phi'=-\int_I\,g\,\phi\ for\ all\ \phi\in C^1_c(I)\}
\end{equation}
and $H^1(I)=W^{1,2}(I)$. 
\end{definition}
在這空間下，我們定義$u'=g$ (這就是弱解)。結合過去學過的泛函列出以下定理。

\begin{proposition} The space $W^{1,p}(I)$ is a Banach space for $1\leq p\leq\infty$. It is reflexive for $1<p<\infty$ and separable for $1\leq p<\infty.$ The space $H^1$ is a separable Hilbert space.
\end{proposition}
\begin{proof} 
See \cite[p.~203]{Bre10}.
\end{proof}
我們列出在這低維度版本的Sobolev embedding theorem。

\begin{theorem}\label{Sobolevembedding} There exists a constant $C$ (depending only on $\lvert I\rvert\leq\infty$) such that 
$$
\lVert u\rVert_{L^\infty(I)}\leq C\lVert u\rVert_{W^{1,p}(I)}\quad\forall u\in W^{1,p}(I),\ \forall 1\leq p\leq\infty.
$$
In other words, $W^{1,p}(I)\subset L^\infty(I)$ with continuous injection for all $1\leq p\leq\infty.$ 
\end{theorem}
\begin{proof} 
See \cite[p.~213]{Bre10}.
\end{proof}
我們解讀這個連續入射是什麼意思，連續性是因為我們在大空間的連續性可以被小空間的連續性控制住，所以連續性是會被小空間所保證的(收斂的部分需要稠密性定理來保證)。至於入射，因為這是從小空間送到大空間，入射性自然可以保證，我們沒有改變函數本身，就一個identity map就好了，這自然是入射，但因為另一個空間比較大，自然沒辦法滿射(就找反例知道哪些函數在大空間卻不在小空間中)。所以這個不等式\eqref{Sineq}主要給我們的是連續性。

\end{document}
